\section{Modelado y evaluación}

\subsection{Modelos considerados}
Se evaluaron cuatro alternativas de clasificación supervisada. La regresión logística se adoptó como \textit{baseline}, por ser una referencia sencilla de interpretar. \textit{Random Forest} y XGBoost permitieron evaluar modelos basados en \textit{árboles}, capaces de representar relaciones no lineales. Por último, se entrenó una red neuronal \textit{Feed-Forward} con dos capas ocultas de 500 neuronas y activación ReLU, \textit{dropout} de 0,2 después de cada capa y salida sigmoide.

La red neuronal se optimizó con Adam y entropía cruzada binaria; se estableció un máximo de 100 \textit{epochs}, lotes de 128 ejemplos y detención temprana con paciencia de cinco \textit{epochs} según el PR-AUC de validación. Los modelos restantes se ejecutaron con los parámetros predefinidos por scikit-learn sin una búsqueda sistemática de hiperparámetros. Cada algoritmo se entrenó en dos variantes: con y sin las dieciocho variables históricas.

\subsection{Métricas y criterio de selección}
La evaluación consideró tres indicadores complementarios. \textit{ROC-AUC} resume la capacidad de ordenar operaciones fraudulentas por encima de las legítimas. \textit{Average Precision} integra el comportamiento de precisión y \textit{recall} a distintos umbrales, y resulta particularmente informativa para un problema de clases desbalanceadas \cite{handbook}. Su expresión es:
\begin{equation}
\label{eq:ap}
AP=\sum_{i}(R_i-R_{i-1})P_i,
\end{equation}
donde $P_i$ y $R_i$ son la precisión y el \textit{recall} de los puntos evaluados.

Para representar la restricción operativa se utilizó CP@$k$. En cada día se agrupan las operaciones por cliente, se conserva el mayor puntaje de riesgo asignado y se seleccionan los $k$ clientes más sospechosos. Un cliente con una o varias operaciones fraudulentas cuenta como una detección correcta. Después de identificarlo, se retira de las evaluaciones de los días posteriores. La métrica general es el promedio de las precisiones diarias:
\begin{equation}
\label{eq:cp}
CP@k=\frac{1}{D}\sum_{d=1}^{D}\frac{C_{d,k}}{k},
\end{equation}
donde $D$ es el número de días y $C_{d,k}$ el número de clientes comprometidos correctamente detectados entre los primeros $k$ del día $d$.

El anteproyecto propuso CP@100 como KPI principal, siguiendo el ejemplo del manual. Sin embargo, en esta simulación el número de clientes comprometidos por día es reducido: en los 23 días de validación hay solamente 142 clientes distintos con fraude. Con una política de retirar clientes ya detectados, el límite de CP@100 es $142/(23\times100)=0,0617$. La regresión logística alcanzó ese valor, lo que hizo poco informativa esta capacidad para distinguir alternativas en la etapa de selección. Por ello se utilizó CP@20 como criterio operativo prioritario, manteniendo AP y ROC-AUC como medidas de apoyo.

\subsection{Resultados en validación}
Los resultados obtenidos con los ocho experimentos se presentan en el Cuadro~\ref{tab:validacion}. Todas las cifras proceden de las salidas calculadas del notebook sobre el conjunto de validación filtrado, correspondiente a los días 127--149.

\begin{table}[!htbp]
\centering
\small
\caption{Comparación de modelos sobre validación}
\label{tab:validacion}
\begin{tabular}{@{}llrrr@{}}
\toprule
\textbf{Modelo} & \textbf{Variables históricas} & \textbf{CP@20} & \textbf{AP} & \textbf{ROC-AUC} \\
\midrule
Regresión logística & Sin & 0,2826 & 0,7555 & 0,9975 \\
Regresión logística & Con & 0,2848 & 0,7985 & 0,9978 \\
Random Forest & Sin & 0,2696 & 0,7091 & 0,9303 \\
Random Forest & Con & 0,2957 & 0,8611 & 0,9959 \\
XGBoost & Sin & 0,2913 & 0,7891 & 0,9984 \\
XGBoost & Con & 0,2870 & 0,7075 & 0,9797 \\
Red neuronal & Sin & 0,2913 & 0,8318 & 0,9984 \\
\textbf{Red neuronal} & \textbf{Con} & \textbf{0,3000} & \textbf{0,8890} & \textbf{0,9993} \\
\bottomrule
\end{tabular}
\par\smallskip
{\footnotesize Valores redondeados a cuatro decimales.}
\end{table}

La ingeniería de variables mejoró la AP en la regresión logística, Random Forest y la red neuronal, aunque no ocurrió lo mismo con XGBoost. En particular, la AP de la red neuronal pasó de 0,8318 sin antecedentes a 0,8890 con ellos, una diferencia de aproximadamente 0,0573 puntos. Su CP@20 aumentó de 0,2913 a 0,3000. En cambio, XGBoost presentó una disminución de AP y CP@20 al incorporar las nuevas variables. Por lo tanto, los resultados respaldan la utilidad de las características históricas para algunos algoritmos, pero no una mejora generalizada.

% \begin{figure}[!htbp]
% \centering
% \begin{tikzpicture}
% \begin{axis}[
%     width=0.91\linewidth,height=6.0cm,
%     ybar,bar width=9pt,
%     ymin=0,ymax=1,
%     ylabel={Average Precision},
%     symbolic x coords={Logística,Forest,XGBoost,Red neuronal},
%     xtick=data,
%     x tick label style={font=\small},
%     y tick label style={font=\small},
%     legend style={at={(0.5,1.14)},anchor=south,legend columns=2,draw=none,font=\small},
%     grid=major,grid style={gray!15},
%     axis x line*=bottom,axis y line*=left,
%     enlarge x limits=0.15
% ]
% \addplot[fill=blue!50,draw=blue!80!black] coordinates {
%  (Logística,0.755506) (Forest,0.709099) (XGBoost,0.789103) (Red neuronal,0.831752)};
% \addplot[fill=pink!50,draw=pink!80!black] coordinates {
%  (Logística,0.798490) (Forest,0.861130) (XGBoost,0.707471) (Red neuronal,0.889045)};
% \legend{Sin variables históricas,Con variables históricas}
% \end{axis}
% \end{tikzpicture}
% \caption{Average Precision con y sin ingeniería de variables, calculada en validación.}
% \label{fig:ap}
% \end{figure}

\begin{figure}[!htbp]
    \centering
    \resizebox{1\textwidth}{!}{%% Creator: Matplotlib, PGF backend
%%
%% To include the figure in your LaTeX document, write
%%   \input{<filename>.pgf}
%%
%% Make sure the required packages are loaded in your preamble
%%   \usepackage{pgf}
%%
%% Also ensure that all the required font packages are loaded; for instance,
%% the lmodern package is sometimes necessary when using math font.
%%   \usepackage{lmodern}
%%
%% Figures using additional raster images can only be included by \input if
%% they are in the same directory as the main LaTeX file. For loading figures
%% from other directories you can use the `import` package
%%   \usepackage{import}
%%
%% and then include the figures with
%%   \import{<path to file>}{<filename>.pgf}
%%
%% Matplotlib used the following preamble
%%   \def\mathdefault#1{#1}
%%   \everymath=\expandafter{\the\everymath\displaystyle}
%%   \IfFileExists{scrextend.sty}{
%%     \usepackage[fontsize=10.000000pt]{scrextend}
%%   }{
%%     \renewcommand{\normalsize}{\fontsize{10.000000}{12.000000}\selectfont}
%%     \normalsize
%%   }
%%   
%%   \ifdefined\pdftexversion\else  % non-pdftex case.
%%     \usepackage{fontspec}
%%     \setmainfont{DejaVuSerif.ttf}[Path=\detokenize{/home/ginoanci/Documentos/Diplomaturas/Trabajos_Practicos/Practicos_Diplo_CD/TF/.venv/lib/python3.13/site-packages/matplotlib/mpl-data/fonts/ttf/}]
%%     \setsansfont{DejaVuSans.ttf}[Path=\detokenize{/home/ginoanci/Documentos/Diplomaturas/Trabajos_Practicos/Practicos_Diplo_CD/TF/.venv/lib/python3.13/site-packages/matplotlib/mpl-data/fonts/ttf/}]
%%     \setmonofont{DejaVuSansMono.ttf}[Path=\detokenize{/home/ginoanci/Documentos/Diplomaturas/Trabajos_Practicos/Practicos_Diplo_CD/TF/.venv/lib/python3.13/site-packages/matplotlib/mpl-data/fonts/ttf/}]
%%   \fi
%%   \makeatletter\@ifpackageloaded{underscore}{}{\usepackage[strings]{underscore}}\makeatother
%%
\begingroup%
\makeatletter%
\begin{pgfpicture}%
\pgfpathrectangle{\pgfpointorigin}{\pgfqpoint{7.900879in}{2.678502in}}%
\pgfusepath{use as bounding box, clip}%
\begin{pgfscope}%
\pgfsetbuttcap%
\pgfsetmiterjoin%
\definecolor{currentfill}{rgb}{1.000000,1.000000,1.000000}%
\pgfsetfillcolor{currentfill}%
\pgfsetlinewidth{0.000000pt}%
\definecolor{currentstroke}{rgb}{1.000000,1.000000,1.000000}%
\pgfsetstrokecolor{currentstroke}%
\pgfsetdash{}{0pt}%
\pgfpathmoveto{\pgfqpoint{0.000000in}{0.000000in}}%
\pgfpathlineto{\pgfqpoint{7.900879in}{0.000000in}}%
\pgfpathlineto{\pgfqpoint{7.900879in}{2.678502in}}%
\pgfpathlineto{\pgfqpoint{0.000000in}{2.678502in}}%
\pgfpathlineto{\pgfqpoint{0.000000in}{0.000000in}}%
\pgfpathclose%
\pgfusepath{fill}%
\end{pgfscope}%
\begin{pgfscope}%
\pgfsetbuttcap%
\pgfsetmiterjoin%
\definecolor{currentfill}{rgb}{1.000000,1.000000,1.000000}%
\pgfsetfillcolor{currentfill}%
\pgfsetlinewidth{0.000000pt}%
\definecolor{currentstroke}{rgb}{0.000000,0.000000,0.000000}%
\pgfsetstrokecolor{currentstroke}%
\pgfsetstrokeopacity{0.000000}%
\pgfsetdash{}{0pt}%
\pgfpathmoveto{\pgfqpoint{0.626435in}{0.341672in}}%
\pgfpathlineto{\pgfqpoint{7.800879in}{0.341672in}}%
\pgfpathlineto{\pgfqpoint{7.800879in}{2.347128in}}%
\pgfpathlineto{\pgfqpoint{0.626435in}{2.347128in}}%
\pgfpathlineto{\pgfqpoint{0.626435in}{0.341672in}}%
\pgfpathclose%
\pgfusepath{fill}%
\end{pgfscope}%
\begin{pgfscope}%
\pgfpathrectangle{\pgfqpoint{0.626435in}{0.341672in}}{\pgfqpoint{7.174444in}{2.005456in}}%
\pgfusepath{clip}%
\pgfsetbuttcap%
\pgfsetmiterjoin%
\definecolor{currentfill}{rgb}{0.560784,0.658824,0.749020}%
\pgfsetfillcolor{currentfill}%
\pgfsetlinewidth{0.803000pt}%
\definecolor{currentstroke}{rgb}{0.435294,0.529412,0.611765}%
\pgfsetstrokecolor{currentstroke}%
\pgfsetdash{}{0pt}%
\pgfpathmoveto{\pgfqpoint{0.952546in}{0.341672in}}%
\pgfpathlineto{\pgfqpoint{1.301951in}{0.341672in}}%
\pgfpathlineto{\pgfqpoint{1.301951in}{1.856806in}}%
\pgfpathlineto{\pgfqpoint{0.952546in}{1.856806in}}%
\pgfpathlineto{\pgfqpoint{0.952546in}{0.341672in}}%
\pgfpathclose%
\pgfusepath{stroke,fill}%
\end{pgfscope}%
\begin{pgfscope}%
\pgfpathrectangle{\pgfqpoint{0.626435in}{0.341672in}}{\pgfqpoint{7.174444in}{2.005456in}}%
\pgfusepath{clip}%
\pgfsetbuttcap%
\pgfsetmiterjoin%
\definecolor{currentfill}{rgb}{0.560784,0.658824,0.749020}%
\pgfsetfillcolor{currentfill}%
\pgfsetlinewidth{0.803000pt}%
\definecolor{currentstroke}{rgb}{0.435294,0.529412,0.611765}%
\pgfsetstrokecolor{currentstroke}%
\pgfsetdash{}{0pt}%
\pgfpathmoveto{\pgfqpoint{2.893684in}{0.341672in}}%
\pgfpathlineto{\pgfqpoint{3.243088in}{0.341672in}}%
\pgfpathlineto{\pgfqpoint{3.243088in}{1.763739in}}%
\pgfpathlineto{\pgfqpoint{2.893684in}{1.763739in}}%
\pgfpathlineto{\pgfqpoint{2.893684in}{0.341672in}}%
\pgfpathclose%
\pgfusepath{stroke,fill}%
\end{pgfscope}%
\begin{pgfscope}%
\pgfpathrectangle{\pgfqpoint{0.626435in}{0.341672in}}{\pgfqpoint{7.174444in}{2.005456in}}%
\pgfusepath{clip}%
\pgfsetbuttcap%
\pgfsetmiterjoin%
\definecolor{currentfill}{rgb}{0.560784,0.658824,0.749020}%
\pgfsetfillcolor{currentfill}%
\pgfsetlinewidth{0.803000pt}%
\definecolor{currentstroke}{rgb}{0.435294,0.529412,0.611765}%
\pgfsetstrokecolor{currentstroke}%
\pgfsetdash{}{0pt}%
\pgfpathmoveto{\pgfqpoint{4.834821in}{0.341672in}}%
\pgfpathlineto{\pgfqpoint{5.184226in}{0.341672in}}%
\pgfpathlineto{\pgfqpoint{5.184226in}{1.924184in}}%
\pgfpathlineto{\pgfqpoint{4.834821in}{1.924184in}}%
\pgfpathlineto{\pgfqpoint{4.834821in}{0.341672in}}%
\pgfpathclose%
\pgfusepath{stroke,fill}%
\end{pgfscope}%
\begin{pgfscope}%
\pgfpathrectangle{\pgfqpoint{0.626435in}{0.341672in}}{\pgfqpoint{7.174444in}{2.005456in}}%
\pgfusepath{clip}%
\pgfsetbuttcap%
\pgfsetmiterjoin%
\definecolor{currentfill}{rgb}{0.560784,0.658824,0.749020}%
\pgfsetfillcolor{currentfill}%
\pgfsetlinewidth{0.803000pt}%
\definecolor{currentstroke}{rgb}{0.435294,0.529412,0.611765}%
\pgfsetstrokecolor{currentstroke}%
\pgfsetdash{}{0pt}%
\pgfpathmoveto{\pgfqpoint{6.775959in}{0.341672in}}%
\pgfpathlineto{\pgfqpoint{7.125364in}{0.341672in}}%
\pgfpathlineto{\pgfqpoint{7.125364in}{2.008116in}}%
\pgfpathlineto{\pgfqpoint{6.775959in}{2.008116in}}%
\pgfpathlineto{\pgfqpoint{6.775959in}{0.341672in}}%
\pgfpathclose%
\pgfusepath{stroke,fill}%
\end{pgfscope}%
\begin{pgfscope}%
\pgfpathrectangle{\pgfqpoint{0.626435in}{0.341672in}}{\pgfqpoint{7.174444in}{2.005456in}}%
\pgfusepath{clip}%
\pgfsetbuttcap%
\pgfsetmiterjoin%
\definecolor{currentfill}{rgb}{0.850980,0.549020,0.584314}%
\pgfsetfillcolor{currentfill}%
\pgfsetlinewidth{0.803000pt}%
\definecolor{currentstroke}{rgb}{0.709804,0.431373,0.470588}%
\pgfsetstrokecolor{currentstroke}%
\pgfsetdash{}{0pt}%
\pgfpathmoveto{\pgfqpoint{1.301951in}{0.341672in}}%
\pgfpathlineto{\pgfqpoint{1.651356in}{0.341672in}}%
\pgfpathlineto{\pgfqpoint{1.651356in}{1.944032in}}%
\pgfpathlineto{\pgfqpoint{1.301951in}{1.944032in}}%
\pgfpathlineto{\pgfqpoint{1.301951in}{0.341672in}}%
\pgfpathclose%
\pgfusepath{stroke,fill}%
\end{pgfscope}%
\begin{pgfscope}%
\pgfpathrectangle{\pgfqpoint{0.626435in}{0.341672in}}{\pgfqpoint{7.174444in}{2.005456in}}%
\pgfusepath{clip}%
\pgfsetbuttcap%
\pgfsetmiterjoin%
\definecolor{currentfill}{rgb}{0.850980,0.549020,0.584314}%
\pgfsetfillcolor{currentfill}%
\pgfsetlinewidth{0.803000pt}%
\definecolor{currentstroke}{rgb}{0.709804,0.431373,0.470588}%
\pgfsetstrokecolor{currentstroke}%
\pgfsetdash{}{0pt}%
\pgfpathmoveto{\pgfqpoint{3.243088in}{0.341672in}}%
\pgfpathlineto{\pgfqpoint{3.592493in}{0.341672in}}%
\pgfpathlineto{\pgfqpoint{3.592493in}{2.068631in}}%
\pgfpathlineto{\pgfqpoint{3.243088in}{2.068631in}}%
\pgfpathlineto{\pgfqpoint{3.243088in}{0.341672in}}%
\pgfpathclose%
\pgfusepath{stroke,fill}%
\end{pgfscope}%
\begin{pgfscope}%
\pgfpathrectangle{\pgfqpoint{0.626435in}{0.341672in}}{\pgfqpoint{7.174444in}{2.005456in}}%
\pgfusepath{clip}%
\pgfsetbuttcap%
\pgfsetmiterjoin%
\definecolor{currentfill}{rgb}{0.850980,0.549020,0.584314}%
\pgfsetfillcolor{currentfill}%
\pgfsetlinewidth{0.803000pt}%
\definecolor{currentstroke}{rgb}{0.709804,0.431373,0.470588}%
\pgfsetstrokecolor{currentstroke}%
\pgfsetdash{}{0pt}%
\pgfpathmoveto{\pgfqpoint{5.184226in}{0.341672in}}%
\pgfpathlineto{\pgfqpoint{5.533631in}{0.341672in}}%
\pgfpathlineto{\pgfqpoint{5.533631in}{1.760474in}}%
\pgfpathlineto{\pgfqpoint{5.184226in}{1.760474in}}%
\pgfpathlineto{\pgfqpoint{5.184226in}{0.341672in}}%
\pgfpathclose%
\pgfusepath{stroke,fill}%
\end{pgfscope}%
\begin{pgfscope}%
\pgfpathrectangle{\pgfqpoint{0.626435in}{0.341672in}}{\pgfqpoint{7.174444in}{2.005456in}}%
\pgfusepath{clip}%
\pgfsetbuttcap%
\pgfsetmiterjoin%
\definecolor{currentfill}{rgb}{0.850980,0.549020,0.584314}%
\pgfsetfillcolor{currentfill}%
\pgfsetlinewidth{0.803000pt}%
\definecolor{currentstroke}{rgb}{0.709804,0.431373,0.470588}%
\pgfsetstrokecolor{currentstroke}%
\pgfsetdash{}{0pt}%
\pgfpathmoveto{\pgfqpoint{7.125364in}{0.341672in}}%
\pgfpathlineto{\pgfqpoint{7.474768in}{0.341672in}}%
\pgfpathlineto{\pgfqpoint{7.474768in}{2.112941in}}%
\pgfpathlineto{\pgfqpoint{7.125364in}{2.112941in}}%
\pgfpathlineto{\pgfqpoint{7.125364in}{0.341672in}}%
\pgfpathclose%
\pgfusepath{stroke,fill}%
\end{pgfscope}%
\begin{pgfscope}%
\pgfsetbuttcap%
\pgfsetroundjoin%
\definecolor{currentfill}{rgb}{0.000000,0.000000,0.000000}%
\pgfsetfillcolor{currentfill}%
\pgfsetlinewidth{0.803000pt}%
\definecolor{currentstroke}{rgb}{0.000000,0.000000,0.000000}%
\pgfsetstrokecolor{currentstroke}%
\pgfsetdash{}{0pt}%
\pgfsys@defobject{currentmarker}{\pgfqpoint{0.000000in}{-0.048611in}}{\pgfqpoint{0.000000in}{0.000000in}}{%
\pgfpathmoveto{\pgfqpoint{0.000000in}{0.000000in}}%
\pgfpathlineto{\pgfqpoint{0.000000in}{-0.048611in}}%
\pgfusepath{stroke,fill}%
}%
\begin{pgfscope}%
\pgfsys@transformshift{1.301951in}{0.341672in}%
\pgfsys@useobject{currentmarker}{}%
\end{pgfscope}%
\end{pgfscope}%
\begin{pgfscope}%
\definecolor{textcolor}{rgb}{0.000000,0.000000,0.000000}%
\pgfsetstrokecolor{textcolor}%
\pgfsetfillcolor{textcolor}%
\pgftext[x=1.301951in,y=0.244450in,,top]{\color{textcolor}{\sffamily\fontsize{10.000000}{12.000000}\selectfont\catcode`\^=\active\def^{\ifmmode\sp\else\^{}\fi}\catcode`\%=\active\def%{\%}Logística}}%
\end{pgfscope}%
\begin{pgfscope}%
\pgfsetbuttcap%
\pgfsetroundjoin%
\definecolor{currentfill}{rgb}{0.000000,0.000000,0.000000}%
\pgfsetfillcolor{currentfill}%
\pgfsetlinewidth{0.803000pt}%
\definecolor{currentstroke}{rgb}{0.000000,0.000000,0.000000}%
\pgfsetstrokecolor{currentstroke}%
\pgfsetdash{}{0pt}%
\pgfsys@defobject{currentmarker}{\pgfqpoint{0.000000in}{-0.048611in}}{\pgfqpoint{0.000000in}{0.000000in}}{%
\pgfpathmoveto{\pgfqpoint{0.000000in}{0.000000in}}%
\pgfpathlineto{\pgfqpoint{0.000000in}{-0.048611in}}%
\pgfusepath{stroke,fill}%
}%
\begin{pgfscope}%
\pgfsys@transformshift{3.243088in}{0.341672in}%
\pgfsys@useobject{currentmarker}{}%
\end{pgfscope}%
\end{pgfscope}%
\begin{pgfscope}%
\definecolor{textcolor}{rgb}{0.000000,0.000000,0.000000}%
\pgfsetstrokecolor{textcolor}%
\pgfsetfillcolor{textcolor}%
\pgftext[x=3.243088in,y=0.244450in,,top]{\color{textcolor}{\sffamily\fontsize{10.000000}{12.000000}\selectfont\catcode`\^=\active\def^{\ifmmode\sp\else\^{}\fi}\catcode`\%=\active\def%{\%}Random Forest}}%
\end{pgfscope}%
\begin{pgfscope}%
\pgfsetbuttcap%
\pgfsetroundjoin%
\definecolor{currentfill}{rgb}{0.000000,0.000000,0.000000}%
\pgfsetfillcolor{currentfill}%
\pgfsetlinewidth{0.803000pt}%
\definecolor{currentstroke}{rgb}{0.000000,0.000000,0.000000}%
\pgfsetstrokecolor{currentstroke}%
\pgfsetdash{}{0pt}%
\pgfsys@defobject{currentmarker}{\pgfqpoint{0.000000in}{-0.048611in}}{\pgfqpoint{0.000000in}{0.000000in}}{%
\pgfpathmoveto{\pgfqpoint{0.000000in}{0.000000in}}%
\pgfpathlineto{\pgfqpoint{0.000000in}{-0.048611in}}%
\pgfusepath{stroke,fill}%
}%
\begin{pgfscope}%
\pgfsys@transformshift{5.184226in}{0.341672in}%
\pgfsys@useobject{currentmarker}{}%
\end{pgfscope}%
\end{pgfscope}%
\begin{pgfscope}%
\definecolor{textcolor}{rgb}{0.000000,0.000000,0.000000}%
\pgfsetstrokecolor{textcolor}%
\pgfsetfillcolor{textcolor}%
\pgftext[x=5.184226in,y=0.244450in,,top]{\color{textcolor}{\sffamily\fontsize{10.000000}{12.000000}\selectfont\catcode`\^=\active\def^{\ifmmode\sp\else\^{}\fi}\catcode`\%=\active\def%{\%}XGBoost}}%
\end{pgfscope}%
\begin{pgfscope}%
\pgfsetbuttcap%
\pgfsetroundjoin%
\definecolor{currentfill}{rgb}{0.000000,0.000000,0.000000}%
\pgfsetfillcolor{currentfill}%
\pgfsetlinewidth{0.803000pt}%
\definecolor{currentstroke}{rgb}{0.000000,0.000000,0.000000}%
\pgfsetstrokecolor{currentstroke}%
\pgfsetdash{}{0pt}%
\pgfsys@defobject{currentmarker}{\pgfqpoint{0.000000in}{-0.048611in}}{\pgfqpoint{0.000000in}{0.000000in}}{%
\pgfpathmoveto{\pgfqpoint{0.000000in}{0.000000in}}%
\pgfpathlineto{\pgfqpoint{0.000000in}{-0.048611in}}%
\pgfusepath{stroke,fill}%
}%
\begin{pgfscope}%
\pgfsys@transformshift{7.125364in}{0.341672in}%
\pgfsys@useobject{currentmarker}{}%
\end{pgfscope}%
\end{pgfscope}%
\begin{pgfscope}%
\definecolor{textcolor}{rgb}{0.000000,0.000000,0.000000}%
\pgfsetstrokecolor{textcolor}%
\pgfsetfillcolor{textcolor}%
\pgftext[x=7.125364in,y=0.244450in,,top]{\color{textcolor}{\sffamily\fontsize{10.000000}{12.000000}\selectfont\catcode`\^=\active\def^{\ifmmode\sp\else\^{}\fi}\catcode`\%=\active\def%{\%}Red neuronal}}%
\end{pgfscope}%
\begin{pgfscope}%
\pgfpathrectangle{\pgfqpoint{0.626435in}{0.341672in}}{\pgfqpoint{7.174444in}{2.005456in}}%
\pgfusepath{clip}%
\pgfsetrectcap%
\pgfsetroundjoin%
\pgfsetlinewidth{0.803000pt}%
\definecolor{currentstroke}{rgb}{0.690196,0.690196,0.690196}%
\pgfsetstrokecolor{currentstroke}%
\pgfsetstrokeopacity{0.200000}%
\pgfsetdash{}{0pt}%
\pgfpathmoveto{\pgfqpoint{0.626435in}{0.341672in}}%
\pgfpathlineto{\pgfqpoint{7.800879in}{0.341672in}}%
\pgfusepath{stroke}%
\end{pgfscope}%
\begin{pgfscope}%
\pgfsetbuttcap%
\pgfsetroundjoin%
\definecolor{currentfill}{rgb}{0.000000,0.000000,0.000000}%
\pgfsetfillcolor{currentfill}%
\pgfsetlinewidth{0.803000pt}%
\definecolor{currentstroke}{rgb}{0.000000,0.000000,0.000000}%
\pgfsetstrokecolor{currentstroke}%
\pgfsetdash{}{0pt}%
\pgfsys@defobject{currentmarker}{\pgfqpoint{-0.048611in}{0.000000in}}{\pgfqpoint{-0.000000in}{0.000000in}}{%
\pgfpathmoveto{\pgfqpoint{-0.000000in}{0.000000in}}%
\pgfpathlineto{\pgfqpoint{-0.048611in}{0.000000in}}%
\pgfusepath{stroke,fill}%
}%
\begin{pgfscope}%
\pgfsys@transformshift{0.626435in}{0.341672in}%
\pgfsys@useobject{currentmarker}{}%
\end{pgfscope}%
\end{pgfscope}%
\begin{pgfscope}%
\definecolor{textcolor}{rgb}{0.000000,0.000000,0.000000}%
\pgfsetstrokecolor{textcolor}%
\pgfsetfillcolor{textcolor}%
\pgftext[x=0.308333in, y=0.288911in, left, base]{\color{textcolor}{\sffamily\fontsize{10.000000}{12.000000}\selectfont\catcode`\^=\active\def^{\ifmmode\sp\else\^{}\fi}\catcode`\%=\active\def%{\%}0.0}}%
\end{pgfscope}%
\begin{pgfscope}%
\pgfpathrectangle{\pgfqpoint{0.626435in}{0.341672in}}{\pgfqpoint{7.174444in}{2.005456in}}%
\pgfusepath{clip}%
\pgfsetrectcap%
\pgfsetroundjoin%
\pgfsetlinewidth{0.803000pt}%
\definecolor{currentstroke}{rgb}{0.690196,0.690196,0.690196}%
\pgfsetstrokecolor{currentstroke}%
\pgfsetstrokeopacity{0.200000}%
\pgfsetdash{}{0pt}%
\pgfpathmoveto{\pgfqpoint{0.626435in}{0.742763in}}%
\pgfpathlineto{\pgfqpoint{7.800879in}{0.742763in}}%
\pgfusepath{stroke}%
\end{pgfscope}%
\begin{pgfscope}%
\pgfsetbuttcap%
\pgfsetroundjoin%
\definecolor{currentfill}{rgb}{0.000000,0.000000,0.000000}%
\pgfsetfillcolor{currentfill}%
\pgfsetlinewidth{0.803000pt}%
\definecolor{currentstroke}{rgb}{0.000000,0.000000,0.000000}%
\pgfsetstrokecolor{currentstroke}%
\pgfsetdash{}{0pt}%
\pgfsys@defobject{currentmarker}{\pgfqpoint{-0.048611in}{0.000000in}}{\pgfqpoint{-0.000000in}{0.000000in}}{%
\pgfpathmoveto{\pgfqpoint{-0.000000in}{0.000000in}}%
\pgfpathlineto{\pgfqpoint{-0.048611in}{0.000000in}}%
\pgfusepath{stroke,fill}%
}%
\begin{pgfscope}%
\pgfsys@transformshift{0.626435in}{0.742763in}%
\pgfsys@useobject{currentmarker}{}%
\end{pgfscope}%
\end{pgfscope}%
\begin{pgfscope}%
\definecolor{textcolor}{rgb}{0.000000,0.000000,0.000000}%
\pgfsetstrokecolor{textcolor}%
\pgfsetfillcolor{textcolor}%
\pgftext[x=0.308333in, y=0.690002in, left, base]{\color{textcolor}{\sffamily\fontsize{10.000000}{12.000000}\selectfont\catcode`\^=\active\def^{\ifmmode\sp\else\^{}\fi}\catcode`\%=\active\def%{\%}0.2}}%
\end{pgfscope}%
\begin{pgfscope}%
\pgfpathrectangle{\pgfqpoint{0.626435in}{0.341672in}}{\pgfqpoint{7.174444in}{2.005456in}}%
\pgfusepath{clip}%
\pgfsetrectcap%
\pgfsetroundjoin%
\pgfsetlinewidth{0.803000pt}%
\definecolor{currentstroke}{rgb}{0.690196,0.690196,0.690196}%
\pgfsetstrokecolor{currentstroke}%
\pgfsetstrokeopacity{0.200000}%
\pgfsetdash{}{0pt}%
\pgfpathmoveto{\pgfqpoint{0.626435in}{1.143855in}}%
\pgfpathlineto{\pgfqpoint{7.800879in}{1.143855in}}%
\pgfusepath{stroke}%
\end{pgfscope}%
\begin{pgfscope}%
\pgfsetbuttcap%
\pgfsetroundjoin%
\definecolor{currentfill}{rgb}{0.000000,0.000000,0.000000}%
\pgfsetfillcolor{currentfill}%
\pgfsetlinewidth{0.803000pt}%
\definecolor{currentstroke}{rgb}{0.000000,0.000000,0.000000}%
\pgfsetstrokecolor{currentstroke}%
\pgfsetdash{}{0pt}%
\pgfsys@defobject{currentmarker}{\pgfqpoint{-0.048611in}{0.000000in}}{\pgfqpoint{-0.000000in}{0.000000in}}{%
\pgfpathmoveto{\pgfqpoint{-0.000000in}{0.000000in}}%
\pgfpathlineto{\pgfqpoint{-0.048611in}{0.000000in}}%
\pgfusepath{stroke,fill}%
}%
\begin{pgfscope}%
\pgfsys@transformshift{0.626435in}{1.143855in}%
\pgfsys@useobject{currentmarker}{}%
\end{pgfscope}%
\end{pgfscope}%
\begin{pgfscope}%
\definecolor{textcolor}{rgb}{0.000000,0.000000,0.000000}%
\pgfsetstrokecolor{textcolor}%
\pgfsetfillcolor{textcolor}%
\pgftext[x=0.308333in, y=1.091093in, left, base]{\color{textcolor}{\sffamily\fontsize{10.000000}{12.000000}\selectfont\catcode`\^=\active\def^{\ifmmode\sp\else\^{}\fi}\catcode`\%=\active\def%{\%}0.4}}%
\end{pgfscope}%
\begin{pgfscope}%
\pgfpathrectangle{\pgfqpoint{0.626435in}{0.341672in}}{\pgfqpoint{7.174444in}{2.005456in}}%
\pgfusepath{clip}%
\pgfsetrectcap%
\pgfsetroundjoin%
\pgfsetlinewidth{0.803000pt}%
\definecolor{currentstroke}{rgb}{0.690196,0.690196,0.690196}%
\pgfsetstrokecolor{currentstroke}%
\pgfsetstrokeopacity{0.200000}%
\pgfsetdash{}{0pt}%
\pgfpathmoveto{\pgfqpoint{0.626435in}{1.544946in}}%
\pgfpathlineto{\pgfqpoint{7.800879in}{1.544946in}}%
\pgfusepath{stroke}%
\end{pgfscope}%
\begin{pgfscope}%
\pgfsetbuttcap%
\pgfsetroundjoin%
\definecolor{currentfill}{rgb}{0.000000,0.000000,0.000000}%
\pgfsetfillcolor{currentfill}%
\pgfsetlinewidth{0.803000pt}%
\definecolor{currentstroke}{rgb}{0.000000,0.000000,0.000000}%
\pgfsetstrokecolor{currentstroke}%
\pgfsetdash{}{0pt}%
\pgfsys@defobject{currentmarker}{\pgfqpoint{-0.048611in}{0.000000in}}{\pgfqpoint{-0.000000in}{0.000000in}}{%
\pgfpathmoveto{\pgfqpoint{-0.000000in}{0.000000in}}%
\pgfpathlineto{\pgfqpoint{-0.048611in}{0.000000in}}%
\pgfusepath{stroke,fill}%
}%
\begin{pgfscope}%
\pgfsys@transformshift{0.626435in}{1.544946in}%
\pgfsys@useobject{currentmarker}{}%
\end{pgfscope}%
\end{pgfscope}%
\begin{pgfscope}%
\definecolor{textcolor}{rgb}{0.000000,0.000000,0.000000}%
\pgfsetstrokecolor{textcolor}%
\pgfsetfillcolor{textcolor}%
\pgftext[x=0.308333in, y=1.492184in, left, base]{\color{textcolor}{\sffamily\fontsize{10.000000}{12.000000}\selectfont\catcode`\^=\active\def^{\ifmmode\sp\else\^{}\fi}\catcode`\%=\active\def%{\%}0.6}}%
\end{pgfscope}%
\begin{pgfscope}%
\pgfpathrectangle{\pgfqpoint{0.626435in}{0.341672in}}{\pgfqpoint{7.174444in}{2.005456in}}%
\pgfusepath{clip}%
\pgfsetrectcap%
\pgfsetroundjoin%
\pgfsetlinewidth{0.803000pt}%
\definecolor{currentstroke}{rgb}{0.690196,0.690196,0.690196}%
\pgfsetstrokecolor{currentstroke}%
\pgfsetstrokeopacity{0.200000}%
\pgfsetdash{}{0pt}%
\pgfpathmoveto{\pgfqpoint{0.626435in}{1.946037in}}%
\pgfpathlineto{\pgfqpoint{7.800879in}{1.946037in}}%
\pgfusepath{stroke}%
\end{pgfscope}%
\begin{pgfscope}%
\pgfsetbuttcap%
\pgfsetroundjoin%
\definecolor{currentfill}{rgb}{0.000000,0.000000,0.000000}%
\pgfsetfillcolor{currentfill}%
\pgfsetlinewidth{0.803000pt}%
\definecolor{currentstroke}{rgb}{0.000000,0.000000,0.000000}%
\pgfsetstrokecolor{currentstroke}%
\pgfsetdash{}{0pt}%
\pgfsys@defobject{currentmarker}{\pgfqpoint{-0.048611in}{0.000000in}}{\pgfqpoint{-0.000000in}{0.000000in}}{%
\pgfpathmoveto{\pgfqpoint{-0.000000in}{0.000000in}}%
\pgfpathlineto{\pgfqpoint{-0.048611in}{0.000000in}}%
\pgfusepath{stroke,fill}%
}%
\begin{pgfscope}%
\pgfsys@transformshift{0.626435in}{1.946037in}%
\pgfsys@useobject{currentmarker}{}%
\end{pgfscope}%
\end{pgfscope}%
\begin{pgfscope}%
\definecolor{textcolor}{rgb}{0.000000,0.000000,0.000000}%
\pgfsetstrokecolor{textcolor}%
\pgfsetfillcolor{textcolor}%
\pgftext[x=0.308333in, y=1.893276in, left, base]{\color{textcolor}{\sffamily\fontsize{10.000000}{12.000000}\selectfont\catcode`\^=\active\def^{\ifmmode\sp\else\^{}\fi}\catcode`\%=\active\def%{\%}0.8}}%
\end{pgfscope}%
\begin{pgfscope}%
\pgfpathrectangle{\pgfqpoint{0.626435in}{0.341672in}}{\pgfqpoint{7.174444in}{2.005456in}}%
\pgfusepath{clip}%
\pgfsetrectcap%
\pgfsetroundjoin%
\pgfsetlinewidth{0.803000pt}%
\definecolor{currentstroke}{rgb}{0.690196,0.690196,0.690196}%
\pgfsetstrokecolor{currentstroke}%
\pgfsetstrokeopacity{0.200000}%
\pgfsetdash{}{0pt}%
\pgfpathmoveto{\pgfqpoint{0.626435in}{2.347128in}}%
\pgfpathlineto{\pgfqpoint{7.800879in}{2.347128in}}%
\pgfusepath{stroke}%
\end{pgfscope}%
\begin{pgfscope}%
\pgfsetbuttcap%
\pgfsetroundjoin%
\definecolor{currentfill}{rgb}{0.000000,0.000000,0.000000}%
\pgfsetfillcolor{currentfill}%
\pgfsetlinewidth{0.803000pt}%
\definecolor{currentstroke}{rgb}{0.000000,0.000000,0.000000}%
\pgfsetstrokecolor{currentstroke}%
\pgfsetdash{}{0pt}%
\pgfsys@defobject{currentmarker}{\pgfqpoint{-0.048611in}{0.000000in}}{\pgfqpoint{-0.000000in}{0.000000in}}{%
\pgfpathmoveto{\pgfqpoint{-0.000000in}{0.000000in}}%
\pgfpathlineto{\pgfqpoint{-0.048611in}{0.000000in}}%
\pgfusepath{stroke,fill}%
}%
\begin{pgfscope}%
\pgfsys@transformshift{0.626435in}{2.347128in}%
\pgfsys@useobject{currentmarker}{}%
\end{pgfscope}%
\end{pgfscope}%
\begin{pgfscope}%
\definecolor{textcolor}{rgb}{0.000000,0.000000,0.000000}%
\pgfsetstrokecolor{textcolor}%
\pgfsetfillcolor{textcolor}%
\pgftext[x=0.308333in, y=2.294367in, left, base]{\color{textcolor}{\sffamily\fontsize{10.000000}{12.000000}\selectfont\catcode`\^=\active\def^{\ifmmode\sp\else\^{}\fi}\catcode`\%=\active\def%{\%}1.0}}%
\end{pgfscope}%
\begin{pgfscope}%
\definecolor{textcolor}{rgb}{0.000000,0.000000,0.000000}%
\pgfsetstrokecolor{textcolor}%
\pgfsetfillcolor{textcolor}%
\pgftext[x=0.252778in,y=1.344400in,,bottom,rotate=90.000000]{\color{textcolor}{\sffamily\fontsize{11.000000}{13.200000}\selectfont\catcode`\^=\active\def^{\ifmmode\sp\else\^{}\fi}\catcode`\%=\active\def%{\%}Average Precision}}%
\end{pgfscope}%
\begin{pgfscope}%
\pgfsetrectcap%
\pgfsetmiterjoin%
\pgfsetlinewidth{0.803000pt}%
\definecolor{currentstroke}{rgb}{0.000000,0.000000,0.000000}%
\pgfsetstrokecolor{currentstroke}%
\pgfsetdash{}{0pt}%
\pgfpathmoveto{\pgfqpoint{0.626435in}{0.341672in}}%
\pgfpathlineto{\pgfqpoint{0.626435in}{2.347128in}}%
\pgfusepath{stroke}%
\end{pgfscope}%
\begin{pgfscope}%
\pgfsetrectcap%
\pgfsetmiterjoin%
\pgfsetlinewidth{0.803000pt}%
\definecolor{currentstroke}{rgb}{0.000000,0.000000,0.000000}%
\pgfsetstrokecolor{currentstroke}%
\pgfsetdash{}{0pt}%
\pgfpathmoveto{\pgfqpoint{7.800879in}{0.341672in}}%
\pgfpathlineto{\pgfqpoint{7.800879in}{2.347128in}}%
\pgfusepath{stroke}%
\end{pgfscope}%
\begin{pgfscope}%
\pgfsetrectcap%
\pgfsetmiterjoin%
\pgfsetlinewidth{0.803000pt}%
\definecolor{currentstroke}{rgb}{0.000000,0.000000,0.000000}%
\pgfsetstrokecolor{currentstroke}%
\pgfsetdash{}{0pt}%
\pgfpathmoveto{\pgfqpoint{0.626435in}{0.341672in}}%
\pgfpathlineto{\pgfqpoint{7.800879in}{0.341672in}}%
\pgfusepath{stroke}%
\end{pgfscope}%
\begin{pgfscope}%
\pgfsetrectcap%
\pgfsetmiterjoin%
\pgfsetlinewidth{0.803000pt}%
\definecolor{currentstroke}{rgb}{0.000000,0.000000,0.000000}%
\pgfsetstrokecolor{currentstroke}%
\pgfsetdash{}{0pt}%
\pgfpathmoveto{\pgfqpoint{0.626435in}{2.347128in}}%
\pgfpathlineto{\pgfqpoint{7.800879in}{2.347128in}}%
\pgfusepath{stroke}%
\end{pgfscope}%
\begin{pgfscope}%
\definecolor{textcolor}{rgb}{0.000000,0.000000,0.000000}%
\pgfsetstrokecolor{textcolor}%
\pgfsetfillcolor{textcolor}%
\pgftext[x=1.127248in,y=1.898473in,,bottom]{\color{textcolor}{\sffamily\fontsize{9.000000}{10.800000}\selectfont\catcode`\^=\active\def^{\ifmmode\sp\else\^{}\fi}\catcode`\%=\active\def%{\%}0.76}}%
\end{pgfscope}%
\begin{pgfscope}%
\definecolor{textcolor}{rgb}{0.000000,0.000000,0.000000}%
\pgfsetstrokecolor{textcolor}%
\pgfsetfillcolor{textcolor}%
\pgftext[x=3.068386in,y=1.805406in,,bottom]{\color{textcolor}{\sffamily\fontsize{9.000000}{10.800000}\selectfont\catcode`\^=\active\def^{\ifmmode\sp\else\^{}\fi}\catcode`\%=\active\def%{\%}0.71}}%
\end{pgfscope}%
\begin{pgfscope}%
\definecolor{textcolor}{rgb}{0.000000,0.000000,0.000000}%
\pgfsetstrokecolor{textcolor}%
\pgfsetfillcolor{textcolor}%
\pgftext[x=5.009524in,y=1.965850in,,bottom]{\color{textcolor}{\sffamily\fontsize{9.000000}{10.800000}\selectfont\catcode`\^=\active\def^{\ifmmode\sp\else\^{}\fi}\catcode`\%=\active\def%{\%}0.79}}%
\end{pgfscope}%
\begin{pgfscope}%
\definecolor{textcolor}{rgb}{0.000000,0.000000,0.000000}%
\pgfsetstrokecolor{textcolor}%
\pgfsetfillcolor{textcolor}%
\pgftext[x=6.950661in,y=2.049783in,,bottom]{\color{textcolor}{\sffamily\fontsize{9.000000}{10.800000}\selectfont\catcode`\^=\active\def^{\ifmmode\sp\else\^{}\fi}\catcode`\%=\active\def%{\%}0.83}}%
\end{pgfscope}%
\begin{pgfscope}%
\definecolor{textcolor}{rgb}{0.000000,0.000000,0.000000}%
\pgfsetstrokecolor{textcolor}%
\pgfsetfillcolor{textcolor}%
\pgftext[x=1.476653in,y=1.985698in,,bottom]{\color{textcolor}{\sffamily\fontsize{9.000000}{10.800000}\selectfont\catcode`\^=\active\def^{\ifmmode\sp\else\^{}\fi}\catcode`\%=\active\def%{\%}0.80}}%
\end{pgfscope}%
\begin{pgfscope}%
\definecolor{textcolor}{rgb}{0.000000,0.000000,0.000000}%
\pgfsetstrokecolor{textcolor}%
\pgfsetfillcolor{textcolor}%
\pgftext[x=3.417791in,y=2.110297in,,bottom]{\color{textcolor}{\sffamily\fontsize{9.000000}{10.800000}\selectfont\catcode`\^=\active\def^{\ifmmode\sp\else\^{}\fi}\catcode`\%=\active\def%{\%}0.86}}%
\end{pgfscope}%
\begin{pgfscope}%
\definecolor{textcolor}{rgb}{0.000000,0.000000,0.000000}%
\pgfsetstrokecolor{textcolor}%
\pgfsetfillcolor{textcolor}%
\pgftext[x=5.358928in,y=1.802141in,,bottom]{\color{textcolor}{\sffamily\fontsize{9.000000}{10.800000}\selectfont\catcode`\^=\active\def^{\ifmmode\sp\else\^{}\fi}\catcode`\%=\active\def%{\%}0.71}}%
\end{pgfscope}%
\begin{pgfscope}%
\definecolor{textcolor}{rgb}{0.000000,0.000000,0.000000}%
\pgfsetstrokecolor{textcolor}%
\pgfsetfillcolor{textcolor}%
\pgftext[x=7.300066in,y=2.154608in,,bottom]{\color{textcolor}{\sffamily\fontsize{9.000000}{10.800000}\selectfont\catcode`\^=\active\def^{\ifmmode\sp\else\^{}\fi}\catcode`\%=\active\def%{\%}0.88}}%
\end{pgfscope}%
\begin{pgfscope}%
\pgfsetbuttcap%
\pgfsetmiterjoin%
\definecolor{currentfill}{rgb}{0.560784,0.658824,0.749020}%
\pgfsetfillcolor{currentfill}%
\pgfsetlinewidth{0.803000pt}%
\definecolor{currentstroke}{rgb}{0.435294,0.529412,0.611765}%
\pgfsetstrokecolor{currentstroke}%
\pgfsetdash{}{0pt}%
\pgfpathmoveto{\pgfqpoint{2.874519in}{2.417424in}}%
\pgfpathlineto{\pgfqpoint{3.152297in}{2.417424in}}%
\pgfpathlineto{\pgfqpoint{3.152297in}{2.514646in}}%
\pgfpathlineto{\pgfqpoint{2.874519in}{2.514646in}}%
\pgfpathlineto{\pgfqpoint{2.874519in}{2.417424in}}%
\pgfpathclose%
\pgfusepath{stroke,fill}%
\end{pgfscope}%
\begin{pgfscope}%
\definecolor{textcolor}{rgb}{0.000000,0.000000,0.000000}%
\pgfsetstrokecolor{textcolor}%
\pgfsetfillcolor{textcolor}%
\pgftext[x=3.263408in,y=2.417424in,left,base]{\color{textcolor}{\sffamily\fontsize{10.000000}{12.000000}\selectfont\catcode`\^=\active\def^{\ifmmode\sp\else\^{}\fi}\catcode`\%=\active\def%{\%}Sin Feature}}%
\end{pgfscope}%
\begin{pgfscope}%
\pgfsetbuttcap%
\pgfsetmiterjoin%
\definecolor{currentfill}{rgb}{0.850980,0.549020,0.584314}%
\pgfsetfillcolor{currentfill}%
\pgfsetlinewidth{0.803000pt}%
\definecolor{currentstroke}{rgb}{0.709804,0.431373,0.470588}%
\pgfsetstrokecolor{currentstroke}%
\pgfsetdash{}{0pt}%
\pgfpathmoveto{\pgfqpoint{4.324945in}{2.417424in}}%
\pgfpathlineto{\pgfqpoint{4.602722in}{2.417424in}}%
\pgfpathlineto{\pgfqpoint{4.602722in}{2.514646in}}%
\pgfpathlineto{\pgfqpoint{4.324945in}{2.514646in}}%
\pgfpathlineto{\pgfqpoint{4.324945in}{2.417424in}}%
\pgfpathclose%
\pgfusepath{stroke,fill}%
\end{pgfscope}%
\begin{pgfscope}%
\definecolor{textcolor}{rgb}{0.000000,0.000000,0.000000}%
\pgfsetstrokecolor{textcolor}%
\pgfsetfillcolor{textcolor}%
\pgftext[x=4.713834in,y=2.417424in,left,base]{\color{textcolor}{\sffamily\fontsize{10.000000}{12.000000}\selectfont\catcode`\^=\active\def^{\ifmmode\sp\else\^{}\fi}\catcode`\%=\active\def%{\%}Con Feature}}%
\end{pgfscope}%
\end{pgfpicture}%
\makeatother%
\endgroup%
}
    \caption{Average Precision con y sin ingeniería de variables, calculada en validación.}
    \label{fig:ap}
\end{figure}


El mejor desempeño de CP@20 y AP en validación correspondió a la red neuronal con ingeniería de variables. Frente a la regresión logística con variables históricas, su CP@20 mejoró 0,0152 puntos y su AP 0,0906 puntos. Sobre esta base se seleccionó la red neuronal para la única evaluación final en prueba.

\subsection{Evaluación final sobre el conjunto de prueba}
El conjunto de prueba filtrado comprendió 58.305 operaciones entre los días 157 y 179, de las cuales 113 fueron fraudulentas. La red neuronal seleccionada produjo los siguientes resultados:

\begin{table}[!htbp]
\centering
\caption{Resultados del modelo seleccionado en prueba}
\label{tab:test}
\begin{tabular}{@{}lrr@{}}
\toprule
\textbf{Métrica} & \textbf{Validación} & \textbf{Prueba} \\
\midrule
Card Precision@20 & 0,3000 & 0,2217 \\
Average Precision & 0,8890 & 0,8902 \\
ROC-AUC & 0,9993 & 0,9993 \\
\bottomrule
\end{tabular}
\par\smallskip
\end{table}

La disminución de CP@20 en prueba debe interpretarse considerando el número de clientes comprometidos disponibles. En validación, un total de 142 clientes distintos con fraude establece un máximo de $142/(23\times20)=0,3087$. El modelo alcanzó $138/(23\times20)=0,3000$. En prueba, los 104 clientes comprometidos fijan un máximo de $104/(23\times20)=0,2261$ y el resultado obtenido equivale a $102/(23\times20)=0,2217$. De este modo, aunque la métrica absoluta es menor en prueba, también lo es el máximo posible bajo el escenario evaluado.

En esta simulación se identificaron 102 de 104 clientes comprometidos presentes en prueba (98,1\%). Esta cifra no equivale a un \textit{recall} de transacciones ni garantiza la misma proporción de detección sobre datos reales. Del mismo modo, los valores elevados de AP y ROC-AUC se interpretan en el contexto del conjunto sintético y no constituyen evidencia suficiente de una generalización a una entidad bancaria.
